\documentclass[11pt,a4paper]{article}
\usepackage[utf8]{inputenc}
\usepackage[T1]{fontenc}
\usepackage[french]{babel}
\usepackage[margin=2.4cm]{geometry}
\usepackage{amsmath,amssymb}
\usepackage{booktabs}
\usepackage{xcolor}
\usepackage[colorlinks=true,linkcolor=blue!50!black,citecolor=blue!50!black]{hyperref}
\usepackage{tcolorbox}
\tcbuselibrary{skins,breakable}
\usepackage{titlesec}
\usepackage{enumitem}
\usepackage{parskip}
\setlength{\parskip}{0.4em}

\definecolor{bleuprincip}{HTML}{1F4E78}
\definecolor{bleuclair}{HTML}{2E74B5}
\definecolor{orangealerte}{HTML}{C55A11}
\definecolor{vertok}{HTML}{548235}

\titleformat{\section}{\Large\bfseries\color{bleuprincip}}{\thesection}{0.6em}{}
\titleformat{\subsection}{\large\bfseries\color{bleuclair}}{\thesubsection}{0.6em}{}

\newtcolorbox{intuition}[1][]{colback=bleuclair!7,colframe=bleuclair,boxrule=0.5pt,arc=2pt,
  left=6pt,right=6pt,top=5pt,bottom=5pt,breakable,title={\small\bfseries Intuition},#1}
\newtcolorbox{definitionbox}[1][]{colback=gray!7,colframe=black!60,boxrule=0.4pt,arc=1pt,
  left=6pt,right=6pt,top=5pt,bottom=5pt,breakable,#1}
\newtcolorbox{resultatbox}[1][]{colback=vertok!8,colframe=vertok,boxrule=1pt,arc=2pt,
  left=7pt,right=7pt,top=5pt,bottom=5pt,breakable,#1}
\newtcolorbox{alertebox}[1][]{colback=orangealerte!7,colframe=orangealerte,boxrule=0.6pt,arc=2pt,
  left=7pt,right=7pt,top=5pt,bottom=5pt,breakable,title={\small\bfseries Correction par rapport à la version précédente},#1}

\newcommand{\cit}[1]{\textsuperscript{\hyperlink{ref#1}{[#1]}}}

\title{\vspace{-1.3cm}\textcolor{bleuprincip}{\Huge\bfseries GDPNow-Maroc}\\[0.3cm]
\textcolor{black}{\Large Repli autorégressif --- Sélection de l'ordre $p$ par AIC}\\[0.15cm]
\textcolor{gray}{\large Méthodologie complète, pourquoi et comment, pour les 7 branches sans lien exploitable}}
\author{}
\date{\today}

\begin{document}
\maketitle
\thispagestyle{empty}

\begin{abstract}
\noindent Ce document remplace une version précédente qui fixait arbitrairement $p=4$ pour les 7 branches sans bridge equation exploitable --- une convention reprise sans justification statistique, exactement le même défaut déjà identifié pour $p=5$ dans le BVAR de la Phase 1. Corrigé ici : chaque branche reçoit son \textbf{propre ordre $p^\star$}, sélectionné par AIC sur une grille $p\in\{1,\dots,8\}$ --- cohérent avec la méthode déjà utilisée pour les bridge equations (sélection de $q$ et $r$).
\end{abstract}

\tableofcontents
\vspace{0.5cm}

% ============================================================================
\section{Pourquoi un modèle purement autorégressif, pour ces 7 branches}
% ============================================================================

\begin{intuition}
Ces 7 branches --- 4 sans aucun indicateur retenu depuis la Phase 2 (Services aux entreprises, Administration publique, Éducation-santé, Autres services), et 3 confirmées sans lien statistique exploitable à l'Étape 4 (Agriculture, Transports, Information-communication) --- n'ont \textbf{aucune bridge equation disponible}. En l'absence de tout indicateur exploitable, la seule information restante pour prévoir leur croissance est \textbf{leur propre passé}.
\end{intuition}

C'est exactement le repli que Higgins\cit{1} utilise pour les sous-composantes sans indicateur mensuel disponible (les dortoirs, dans son propre travail) --- une méthode reconnue, pas un pis-aller improvisé.

\begin{equation}
\Delta\log(VA_t) = \alpha + \sum_{k=1}^{p}\gamma_k \Delta\log(VA_{t-k}) + \varepsilon_t
\label{eq:ar_p}
\end{equation}

% ============================================================================
\section{Pourquoi l'ordre $p$ doit être choisi, pas fixé d'avance}
% ============================================================================

\subsection{L'erreur de la version précédente}

\begin{alertebox}
La version précédente fixait $p=4$ pour toutes les branches, justifié uniquement par la convention « un retard par trimestre de l'année écoulée ». \textbf{Exactement le même défaut} que celui déjà documenté pour $p=5$ dans le BVAR de la Phase 1 --- une valeur copiée, jamais testée contre les données réelles de chaque branche.
\end{alertebox}

\subsection{Pourquoi un $p$ différent par branche ne pose aucun problème ici}

\begin{intuition}
Dans le BVAR de la Phase 1, un $p$ commun était \textbf{structurellement nécessaire} --- toutes les branches y forment un seul système à 16 équations simultanées (une matrice compagne commune), incompatible avec des retards différents par équation. \textbf{Ici, ce n'est plus le cas} : chaque branche a son propre modèle AR \emph{indépendant}, sans aucune contrainte partagée. Rien n'empêche mathématiquement Agriculture d'avoir $p=4$ pendant que Transports a $p=8$ --- exactement comme $q$ et $r$ variaient déjà librement d'une bridge equation à l'autre (Étape 4), sans qu'aucun $q$ ou $r$ commun n'ait jamais été imposé.
\end{intuition}

% ============================================================================
\section{Comment --- la sélection AIC, pas à pas}
% ============================================================================

\subsection{La grille testée}

Pour chaque branche, tous les ordres $p\in\{1,\dots,8\}$ sont estimés par MCO :
\begin{equation}
p^\star = \operatorname*{arg\,min}_{p\in\{1,\dots,8\}} \text{AIC}(p), \qquad \text{AIC} = 2k - 2\ln(\hat{L})
\label{eq:aic_ar}
\end{equation}

\subsection{La contrainte d'échantillon commun}

\begin{definitionbox}
Les 8 modèles (un par valeur de $p$) sont d'abord tous estimés sur le \textbf{même échantillon} --- celui compatible avec le plus grand ordre testé ($p=8$, qui perd le plus d'observations en début de période à cause des retards). \textbf{Indispensable} pour que les 8 AIC soient directement comparables entre eux : comparer un AIC calculé sur 40 observations à un AIC calculé sur 44 n'aurait aucun sens, l'AIC dépendant de la taille de l'échantillon.
\end{definitionbox}

\subsection{Réestimation finale}

Une fois $p^\star$ connu, le modèle final est réestimé sur \textbf{toutes} les observations disponibles pour cet ordre précis (pas seulement l'échantillon commun à $p=8$) --- pour ne pas perdre inutilement des points une fois le bon ordre déterminé.

\subsection{Prévision}

À partir des $p^\star$ derniers $\Delta\log$ connus (sur toute la série disponible, pas seulement le train), une prévision hors-échantillon est produite pour le trimestre suivant.

% ============================================================================
\section{Résultats}
% ============================================================================

\begin{resultatbox}
\textbf{7 branches sur 7 estimées avec succès.} Sur les 7, \textbf{4 confirment $p^\star=4$} (la convention initiale se révèle correcte pour elles) --- mais \textbf{2 branches (Transports, Autres services) ont un ordre optimal nettement plus élevé ($p^\star=8$)}, avec un gain de pouvoir explicatif substantiel une fois le bon ordre trouvé.
\end{resultatbox}

\begin{center}
\small
\begin{tabular}{lrrrr}
\toprule
\textbf{Branche} & \textbf{$p^\star$} & \textbf{$R^2$} & \textbf{$R^2$ (ancien, $p=4$)} & \textbf{Gain} \\
\midrule
Autres services & \textbf{8} & 0,804 & 0,556 & $+0{,}248$ \\
Transports & \textbf{8} & 0,482 & 0,231 & $+0{,}251$ \\
Administration publique & 4 & 0,450 & 0,450 & $0$ \\
Éducation-santé & 4 & 0,424 & 0,424 & $0$ \\
Services aux entreprises & 4 & 0,246 & 0,246 & $0$ \\
Agriculture & 4 & 0,186 & 0,186 & $0$ \\
Information-communication & 3 & 0,093 & 0,105 & $-0{,}012$ \\
\bottomrule
\end{tabular}
\end{center}

\begin{intuition}
Information-communication est le seul cas où le $R^2$ \emph{baisse} légèrement en changeant d'ordre ($p^\star=3$ contre $p=4$ imposé avant) --- rappel utile que l'AIC ne maximise pas le $R^2$ en soi, il pénalise la complexité : un modèle plus simple, même avec un $R^2$ marginalement inférieur, peut être préféré s'il généralise mieux.
\end{intuition}

\subsection{Exemple détaillé --- Autres services (grille AIC complète)}

\begin{center}
\small
\begin{tabular}{lrrrrrrrr}
\toprule
$p$ & 1 & 2 & 3 & 4 & 5 & 6 & 7 & \textbf{8} \\
\midrule
AIC & $-200{,}1$ & $-198{,}1$ & $-197{,}8$ & $-216{,}0$ & $-237{,}7$ & $-239{,}2$ & $-238{,}4$ & $\mathbf{-244{,}9}$ \\
\bottomrule
\end{tabular}
\end{center}

\begin{intuition}
Un vrai minimum global à $p=8$, pas un artefact isolé --- amélioration progressive et quasi monotone à partir de $p=4$, confirmant que cette branche a une dynamique dépendant de plusieurs trimestres passés (possiblement un cycle annuel complet, voire plus).
\end{intuition}

% ============================================================================
\section{Sorties produites}
% ============================================================================

\begin{resultatbox}
Pour chaque branche : \texttt{resultats/<branche>\_AR\_coefficients.csv} (le modèle final retenu), \texttt{resultats/<branche>\_AR\_grille\_AIC.csv} (les 8 AIC testés, pour vérification), \texttt{figures/<branche>\_AR\_ajustement.png}. Bilan global dans \texttt{resultats/recapitulatif\_AR.csv}.
\end{resultatbox}

\section*{Références}
\addcontentsline{toc}{section}{Références}
\begin{enumerate}[leftmargin=1.6em, label=\textbf{[\arabic*]}]
\item \hypertarget{ref1}{} Higgins, P. (2014). \textit{GDPNow: A Model for GDP ``Nowcasting''}. Federal Reserve Bank of Atlanta, Working Paper 2014-7.
\end{enumerate}

\end{document}
